\documentclass[11pt]{article}
\usepackage[margin=1in]{geometry}
\usepackage{amsmath,amssymb,amsthm,booktabs,graphicx,microtype}
\usepackage[T1]{fontenc}
\usepackage{xurl}
\usepackage[colorlinks=true,linkcolor=blue,citecolor=blue,urlcolor=blue]{hyperref}
\newtheorem{proposition}{Proposition}
\setlength{\emergencystretch}{3em}
\title{Two Negative Results in Post-Training Quantization:\\
Polar Weight Codes and Self-Consistent Calibration}
\author{Working research draft}
\date{September 21, 2026}
\begin{document}
\maketitle
\begin{center}\small\textbf{Status: working draft. Not submission ready.}\\
No quantum-hardware, quantum-advantage, or quantization-aware-training result is reported.
\end{center}

\begin{abstract}
We report two hypotheses about four-bit weight quantization, both of which we
pursued at length and both of which we now refute, together with the small set
of findings that survive. The first is that a quantum-inspired polar code --
storing a magnitude and a phase index per pair of real weights -- preserves
language-model loss better than an equal-byte real 4-bit code. It does, but only
against a control with a uniform codebook. Giving the scalar control the same
per-tensor Lloyd--Max fitting reverses the ranking on every tensor and on the
complete validation stream, and measuring an unconstrained 256-point
two-dimensional codebook shows the polar arrangement is structurally capped
$6$--$8\%$ above the fitted scalar code while a free 2-D codebook sits
$18$--$24\%$ below it. The second hypothesis is that activation-aware
calibration should be solved as a self-consistent fixed point, by analogy with
the Kohn--Sham loop in density functional theory. It should not: the
forward-only dependency graph is acyclic, so one sequential pass is exact, and
our iteration measurably underperforms that pass. A bidirectional reformulation
using downstream loss sensitivity does identify a genuinely cyclic coupling
$8.41\times$ larger than the one the literature corrects, but the metric it
implies has no rank skill across layers, and its single correct call gives no
clear advantage over spending the same bytes uniformly. Both hypotheses failed the
same way -- they beat a control that had been given less tuning than the method
-- and we document that pattern as the paper's main methodological contribution.
\end{abstract}

\section{Introduction}

This paper is a report of two failures and a short list of what we learned from
them. We state that plainly because intermediate drafts of this work claimed
both hypotheses as positive results, and the corrections matter more than the
original claims did.

\paragraph{Hypothesis 1: polar weight codes.} Pair two real weights as
$z=w_a+iw_b$ and store a three-bit magnitude and a five-bit phase index, the
same eight index bits two real 4-bit weights would take. An earlier draft
reported that this preserved frozen Qwen3.5-0.8B validation loss better than an
equal-byte real 4-bit control across all eighteen DeltaNet input projections.
Section~\ref{sec:part1} shows the advantage came from the control, not the code.

\paragraph{Hypothesis 2: self-consistent calibration.} Activation-aware
quantizers calibrate against activations measured on the full-precision model
but are applied to a model in which many layers have been replaced. Writing the
missing condition as a fixed point $Q^\star=\mathcal Q(H[Q^\star])$ suggests the
Kohn--Sham self-consistent field loop as a solution method.
Section~\ref{sec:part2} shows the forward-only version of this problem is
acyclic and therefore solved exactly by one sequential pass -- as two published
closed-form methods already exploit -- and that a bidirectional reformulation,
while identifying a real and much larger coupling, does not convert into a
usable method.

\paragraph{The common failure mode.} Both hypotheses survived as long as they
did because each was compared against a control that had received less tuning
than the proposed method. The polar code was matched on bytes but not on
codebook adaptation. The layer-protection scheme was compared against another
selection criterion rather than against uniformly spending the same bytes. In
both cases the effect disappeared the moment the control was brought to equal
strength. Section~\ref{sec:discussion} develops this, because we think it is the
most transferable thing in the paper.

\paragraph{What survives.} Three findings, none of them the original hypothesis,
are reported in Section~\ref{sec:survives}: uniformly spending a few more bits
beats every layer-selective scheme we tested; the single most damaging layer to
quantize is the one that weight MSE ranks as the safest, which is a concrete
mechanism for a failure this project observed repeatedly without explaining; and
paired 2-D vector quantization reaches parity with a well-tuned scalar code
where the polar code does not.

\paragraph{What this is not.} Polar notation, complex-valued networks, and
rotation-block weights are prior art, and a complex multiplication by
$\rho e^{i\phi}$ is exactly the real $2\times2$ rotation block. Nothing here is
a quantum algorithm, uses quantum hardware, or demonstrates quantum advantage.
The DFT correspondence in Section~\ref{sec:part2} is a numerical-methods analogy
about fixed-point iteration, and Section~\ref{sec:acyclic} reports that it is
the wrong analogy. We report storage in counted bytes and quality in held-out
loss; we do \emph{not} report an inference speedup.

\section{Related Work and Claim Boundaries}

\paragraph{Post-training quantization.} GPTQ uses approximate second-order
information to compensate rounding error within a block \cite{gptq}; AWQ scales
channels by activation statistics before rounding \cite{awq}. GPTAQ's asymmetric
calibration \cite{gptaq} and CoreQ's mismatch correction \cite{coreq} both
address accumulated error from previously quantized layers in closed form and in
a single pass; they are the prior art that makes our Hypothesis 2 redundant in
its forward-only form, and we found them only after building the iteration. Our
GPTQ- and AWQ-style baselines are local implementations, because available
production backends do not expose Qwen3.5's custom DeltaNet projections.

\paragraph{Complex and phase representations.} Complex-valued networks and
phase-encoded weights predate this work \cite{phaseweights,deepcomplex}, and the
Phase-Coherent Transformer studies phase-preserving attention \cite{pct}.
Expressing a weight as $qe^{i\phi}$ was never a claimed novelty, and
Section~\ref{sec:ceiling} reports it is not even a good code.

\paragraph{Recurrent structure and its quantization.} Unitary recurrent networks
use structured norm-preserving transformations \cite{urnn}; S4D uses complex
diagonal state dynamics \cite{s4d}; QS4D studies quantization-aware training of
structured state-space models \cite{qs4d}; Quamba2 develops post-training
quantization for Mamba backbones \cite{quamba2}. Gated DeltaNet supplies the
recurrent half of our target checkpoints \cite{gdn}, and concurrent work reports
that it tolerates 4-bit quantization well \cite{gdnquant}.

\paragraph{Vector quantization.} Joint quantization of weight groups with a
learned codebook is established practice, and our 2-D codec is a deliberately
simple instance. Our contribution there is the measured ceiling it approaches,
not the codec.

\section{Experimental Protocol}
\label{sec:protocol}

All language-model experiments use the official 752M-parameter
Qwen3.5-0.8B-Base checkpoint as a frozen pretrained model and replace all
eighteen \texttt{linear\_attn.in\_proj\_qkv.weight} tensors. All remaining
weights stay BF16, so reported byte reductions describe that tensor subset and
not whole-checkpoint compression.

\paragraph{Budget.} Every codec spends four index bits per real weight and one
FP32 scale per output row. Per-tensor codebook headers are counted and stated:
16 FP32 for the fitted scalar code, 8 FP32 plus one offset for polar, and
$256\times2$ FP32 (2\,KB) for the 2-D code. Against these $6144\times1024$
matrices the largest header is $2.6\times10^{-3}$ bits per weight, or $0.065\%$
of the four-bit index budget, so the comparison is byte-matched to within
rounding. Where an arm spends more (Table~\ref{tab:protect}), the payload ratio
is stated and matched controls are supplied at the same ratio.

\paragraph{Evaluation.} Unless stated otherwise we evaluate every one of the
2{,}041 contiguous 128-target blocks of the Qwen-tokenized WikiText-2 validation
stream (261{,}248 targets), resetting context at each block boundary. This is
the all-token protocol; an earlier phase established that a
last-token-per-example protocol can reverse codec rankings, so we fix one
objective. Confidence intervals are paired percentile bootstraps over validation
blocks with 5{,}000 replicates. Screening experiments use reduced slices, stated
per table.

\paragraph{Calibration.} Windows of 129 tokens are drawn from the WikiText-2
\emph{training} split with fixed seeds; no validation target is used for any
selection. Where codecs make a per-layer choice, every codec receives the same
selection rule.

\paragraph{Verification.} Every codec is checked to decode exactly from what it
claims to store -- indices, row scales, and the per-tensor codebook header --
and the polar nearest-neighbour claim is verified in float64 against
brute-force search over all 256 codepoints. Byte accounting is asserted against
the real-4 budget.

\section{Part I: Paired Weight Codes}
\label{sec:part1}

\subsection{The hypothesis and the control}

The codec pairs two real weights as $z=w_a+iw_b$ and stores a three-bit
magnitude index and a five-bit phase index per pair -- the same eight index bits
two real 4-bit weights would take. An earlier draft of this work reported that
the code preserved frozen Qwen3.5-0.8B validation loss better than an
equal-byte real 4-bit control across all eighteen DeltaNet input projections,
with $\Delta$NLL $+0.009333$ against the control's $+0.021945$.

That comparison used a real 4-bit control whose codebook is a \emph{uniform}
16-level grid on $[-1,1]$. The polar code's effective codebook is not uniform: a
radius grid crossed with a phase grid produces annular sectors, so the polar
code was receiving a shape adaptation the control was denied. The asymmetry is
in the tuning, not in the representation.

We therefore add a matched control: the same scalar code with a per-tensor
16-level Lloyd--Max codebook fitted to that tensor's own normalized weight
distribution, and the identical per-row FP32 scale. The extra storage is 16 FP32
per tensor, under $10^{-4}$ bits per weight for these matrices.

\subsection{A ceiling for the whole family}
\label{sec:ceiling}

Before asking whether polar helps, we ask how much room exists at eight index
bits per pair, and whether the polar arrangement uses it well. We fit an
\emph{unconstrained} 256-point two-dimensional codebook to the pairs by
$k$-means. Any scheme that pairs two real weights and spends eight index bits on
the pair -- polar, Cartesian, rotated, or otherwise -- is a constrained special
case, so its distortion lower-bounds the entire family at this budget.

\begin{table}[ht]
\centering\small
\setlength{\tabcolsep}{4pt}
\caption{Weight reconstruction at eight index bits per pair, relative to the Lloyd-fitted
real-4 control (values below $1.000$ beat it). \texttt{vq2d} is an unconstrained 256-point
2-D codebook and lower-bounds distortion for every paired code at this budget.
Every entry uses the best of four pairings and a per-row FP32 scale.}
\label{tab:ceiling}
\begin{tabular}{llrrrrr}\toprule
Checkpoint & Tensor & real4-Lloyd MSE & real4-unif. & polar-legacy & polar-fitted & vq2d\\\midrule
Qwen3.5-0.8B & DeltaNet QKV L0 & 3.252e-06 & 1.308 & 1.194 & 1.065 & \textbf{0.759}\\
Qwen3.5-0.8B & DeltaNet QKV L8 & 3.135e-06 & 1.214 & 1.156 & 1.064 & \textbf{0.793}\\
Qwen3.5-0.8B & DeltaNet QKV L16 & 5.039e-06 & 1.240 & 1.165 & 1.059 & \textbf{0.795}\\
Qwen3.8-27B & DeltaNet QKV L0 & 2.563e-06 & 1.186 & 1.162 & 1.073 & \textbf{0.818}\\
Qwen3.8-27B & Attention Q L3 & 3.317e-06 & 1.217 & 1.154 & 1.060 & \textbf{0.791}\\
Qwen3.8-27B & MLP up L3 & 1.182e-06 & 1.169 & 1.163 & 1.081 & \textbf{0.817}\\
\bottomrule
\end{tabular}
\end{table}


Three facts follow from Table~\ref{tab:ceiling}, on every tensor measured across
two checkpoints.

\begin{enumerate}
\item \textbf{The reported advantage is an artifact of the control.} Polar beats
the uniform scalar code, as reported, but the Lloyd-fitted scalar code beats
polar on all six tensors.
\item \textbf{The polar family is structurally capped.} Even after we repair its
assignment rule, fit its radius grid, and sweep every magnitude/phase bit split
(Section~\ref{sec:polarfix}), the best polar code remains $6$--$8\%$ above the
fitted scalar code.
\item \textbf{The pairing idea is sound.} A free 2-D codebook sits $18$--$24\%$
\emph{below} the fitted scalar code at the same index budget. The headroom is
real; the polar parameterization fails to reach it.
\end{enumerate}

\subsection{Why repairing the polar assignment rule changes almost nothing}
\label{sec:polarfix}

The original quantizer rounds magnitude and phase independently, which is not
nearest-neighbour assignment in the $8\times32$ codebook. For a decoded phase
ray $\theta_p$ and a source $z=re^{i\theta}$,
\begin{equation}
 |z-g_m e^{i\theta_p}|^2 = r^2\sin^2\delta + (g_m - r\cos\delta)^2,
 \qquad \delta=\theta-\theta_p ,
\end{equation}
so the optimal radius level is the one nearest the \emph{projection}
$r\cos\delta$, not the one nearest $r$. We implemented exact nearest-neighbour
assignment over all 256 codepoints, verified in float64 against brute force
(maximum relative slack $2.9\times10^{-7}$, zero of 16{,}384 pairs suboptimal).
It reduced weight MSE by $0.035\%$.

The reason is quantitative and retrospectively explains a long series of null
results. With five phase bits the nearest-ray error satisfies
$|\delta|\le\pi/32$, hence $1-\cos\delta\le4.8\times10^{-3}$. The projection
correction is at most half a percent of the radius, while one step of the
three-bit radius grid is about $14\%$ of the row scale, so the correction almost
never changes which radius level wins.

Every phase-side refinement attempted across this project -- learned global
phase offsets, activation-weighted phase choice, output-error phase selection --
was operating on a quantity bounded to be small, which is why each produced
sub-noise, sign-unstable effects across layers and corpora. \textbf{The binding
constraint is the radius grid, not the phase grid.} The bit-split sweep confirms
it: $8\times32$ is optimal on every tensor, and moving bits toward the phase is
catastrophic ($4\times64$ roughly $3\times$ worse, $2\times128$ roughly
$14\times$ worse).

\subsection{Held-out loss, across three selection protocols}

Weight reconstruction does not settle the question; this project has repeatedly
found that weight-error rankings fail to predict task rankings. We therefore
replaced all eighteen projections simultaneously and evaluated the complete
WikiText-2 validation stream. Because the earlier draft selected each layer's
pairing by calibration NLL rather than by weighted error, and because that is a
stronger rule, we report both, plus a reproduction of the earlier draft's exact
calibration size.

\begin{table}[ht]
\centering\small
\caption{Frozen Qwen3.5-0.8B with all 18 DeltaNet input projections quantized
simultaneously, evaluated on all 2041 contiguous 128-target
WikiText-2 validation blocks. Every codec spends four index bits per real weight
and one FP32 row scale. Intervals are paired block bootstraps against BF16
(5{,}000 replicates). Lower is better.
$^{a}$the earlier draft's control; $^{b}$the matched control;
$^{c}$the earlier draft's method.}
\label{tab:main}
\begin{tabular}{lrrrl}\toprule
Method & NLL & Perplexity & $\Delta$NLL & 95\% CI\\\midrule
BF16 reference & 3.436670 & 31.083 & --- & ---\\
\midrule
real-4, uniform codebook$^{a}$ & 3.458469 & 31.768 & +0.021798 & $[+0.020841,\,+0.022752]$\\
real-4, Lloyd codebook$^{b}$ & 3.451102 & 31.535 & +0.014431 & $[+0.013626,\,+0.015235]$\\
\quad + sequential calibration & 3.451124 & 31.536 & +0.014454 & $[+0.013641,\,+0.015255]$\\
polar 3+5, decoupled rounding$^{c}$ & 3.456446 & 31.704 & +0.019776 & $[+0.018821,\,+0.020732]$\\
\quad + sequential calibration & 3.456446 & 31.704 & +0.019776 & $[+0.018821,\,+0.020732]$\\
polar 3+5, exact NN + fitted radii & 3.452258 & 31.572 & +0.015588 & $[+0.014736,\,+0.016421]$\\
\quad + sequential calibration & 3.451662 & 31.553 & +0.014992 & $[+0.014156,\,+0.015820]$\\
paired 2-D codebook (256 points) & 3.445063 & 31.345 & +0.008393 & $[+0.007720,\,+0.009069]$\\
\quad + sequential calibration & 3.445670 & 31.364 & +0.009000 & $[+0.008321,\,+0.009679]$\\
\bottomrule
\end{tabular}
\end{table}

\begin{table}[ht]
\centering\small
\caption{Paired head-to-head differences on the same validation blocks.
A negative value favours the first-named method.}
\label{tab:head}
\begin{tabular}{lrll}\toprule
Comparison & $\Delta$NLL & 95\% CI & Interval\\\midrule
polar (earlier) $-$ real-4 Lloyd & +0.005345 & $[+0.004282,\,+0.006407]$ & excludes 0\\
polar (repaired) $-$ real-4 Lloyd & +0.001156 & $[+0.000156,\,+0.002166]$ & excludes 0\\
paired 2-D $-$ real-4 Lloyd & -0.006038 & $[-0.006952,\,-0.005104]$ & excludes 0\\
paired 2-D $-$ polar (earlier) & -0.011383 & $[-0.012410,\,-0.010380]$ & excludes 0\\
\bottomrule
\end{tabular}
\end{table}


\begin{table}[ht]
\centering\small
\setlength{\tabcolsep}{4pt}
\caption{One protocol for every codec: per-layer candidate selection by calibration NLL
against an otherwise-BF16 model, 16 training windows, then cumulative application.
Codecs that can use activation statistics receive the identical BF16-measured statistic.
Complete WikiText-2 validation stream; paired block bootstraps against BF16.}
\label{tab:fair}
\begin{tabular}{lrrl}\toprule
Method & NLL & $\Delta$NLL & 95\% CI\\\midrule
real-4, uniform codebook (earlier draft's control) & 3.460834 & +0.024164 & $[+0.023088,\,+0.025239]$\\
polar 3+5 (earlier draft's method) & 3.452573 & +0.015902 & $[+0.015006,\,+0.016825]$\\
real-4, Lloyd codebook (matched control) & 3.445209 & +0.008539 & $[+0.007662,\,+0.009403]$\\
paired 2-D codebook & 3.445815 & +0.009145 & $[+0.008325,\,+0.009950]$\\
\textbf{paired 2-D codebook, activation-weighted} & 3.444363 & +0.007693 & $[+0.007029,\,+0.008398]$\\
\midrule
\multicolumn{4}{l}{\footnotesize Earlier draft's published polar figure: $+0.009333$.}\\
\bottomrule
\end{tabular}
\end{table}

\begin{table}[ht]
\centering\small
\setlength{\tabcolsep}{4pt}
\caption{Paired differences under the single protocol of Table~\ref{tab:fair}.
A negative value favours the first-named method.}
\label{tab:fairhead}
\begin{tabular}{lrll}\toprule
Comparison & $\Delta$NLL & 95\% CI & Interval\\\midrule
polar $-$ matched scalar control & +0.007364 & $[+0.006332,\,+0.008376]$ & excludes 0\\
polar $-$ paired 2-D (act.-weighted) & +0.008210 & $[+0.007205,\,+0.009186]$ & excludes 0\\
paired 2-D $-$ same, act.-weighted & +0.001452 & $[+0.000556,\,+0.002327]$ & excludes 0\\
matched scalar $-$ paired 2-D (act.-weighted) & +0.000846 & $[-0.000169,\,+0.001822]$ & \textbf{includes 0}\\
\bottomrule
\end{tabular}
\end{table}


The ranking is stable wherever the protocol can discriminate:

\begin{center}\small
\begin{tabular}{lrrl}\toprule
Protocol & polar & real-4 Lloyd & Discriminates?\\\midrule
Weighted-error selection, 16 windows & $+0.019776$ & $+0.014717$ & yes, Lloyd\\
NLL selection, 16 windows & $+0.015902$ & $+0.008539$ & yes, Lloyd\\
NLL selection, 4 windows, 3 seeds & $0.010945\pm0.000422$ & $0.010554\pm0.003349$ & \textbf{no}\\\bottomrule
\end{tabular}
\end{center}

The third row is the earlier draft's own calibration size, and we report that it
does \emph{not} decide the question. Over three seeds the two codecs are
statistically tied and on one seed polar wins outright. The matched control's
seed-to-seed spread is eight times larger, because its clip-factor search is a
finer decision than polar's four-way pairing choice and 512 calibration tokens
cannot resolve it. That is a narrow but genuine point in polar's favour: coarser
decisions are more stable when calibration data is scarce.

What the small-calibration sweep does establish is that all three of our polar
reproductions fall in $[+0.010515,+0.011358]$, so the earlier draft's published
$+0.009333$ lies outside the spread we could reproduce and is best read as a
fortunate draw from a high-variance selector.

\subsection{What survives: parity, not superiority}

Under a single protocol in which every codec receives per-layer NLL selection
and the codecs that can use activation statistics receive the identical
statistic, the activation-weighted 2-D codebook gives the best point estimate of
anything we tested. Two of its three comparisons are conclusive and one is not:

\begin{itemize}
\item It beats polar decisively, $+0.008210$ $[+0.007205,+0.009186]$.
\item Activation weighting is what makes the difference: the unweighted 2-D
 codebook is worse by $+0.001452$ $[+0.000556,+0.002327]$.
\item Against the matched scalar control the paired difference is $+0.000846$
 with interval $[-0.000169,+0.001822]$, which \emph{includes zero}.
\end{itemize}

We therefore claim parity, not superiority: the 2-D codec reaches the
well-tuned scalar code's regime and may exceed it, while the polar code is
significantly behind both. Replacing the polar grid with a learned codebook also
removes the obstacle that defeated our packed kernels -- decoding a polar index
requires a sine and a cosine per pair, which made decode-plus-GEMM
$3$--$4.5\times$ slower than resident BF16, whereas decoding a 2-D codebook
index is a 256-entry lookup that fits in shared memory. We have not written or
benchmarked that kernel, so this removes a known blocker rather than
demonstrating a speedup.

\subsection{Summary of Part I}

Polar magnitude--phase weight coding loses to a real 4-bit code once the scalar
code receives the same per-tensor codebook fitting, and it is structurally
capped several percent above that control regardless of how its assignment rule,
radius grid, and bit split are tuned. The measurement that establishes this also
locates the headroom: an unconstrained 2-D codebook over the same pairs sits far
below the scalar code at an identical index budget. The pairing idea was right
and the polar parameterization was wrong -- but the resulting codec reaches
parity with a well-tuned scalar code rather than beating it.


\section{Part II: Self-Consistent Calibration}
\label{sec:part2}

\subsection{The calibration fixed point}

Activation-aware post-training quantization chooses a code for layer $\ell$ by
minimizing a metric weighted by that layer's input statistics
$H^{xx}_\ell=\mathbb E[x_\ell x_\ell^\top]$, estimated on calibration data.
GPTQ and AWQ measure those statistics on the full-precision model
\cite{gptq,awq}. Once many layers are replaced, the deployed model presents
layer $\ell$ with a different distribution from the one its codebook was fitted
to, so the conventional procedure computes
$Q_{\text{one-shot}}=\mathcal Q(H[W_{\rm BF16}])$ while the quantity describing
the deployed model is the fixed point
\begin{equation}
 Q^\star=\mathcal Q\bigl(H[Q^\star]\bigr).
 \label{eq:fp}
\end{equation}

Equation~\eqref{eq:fp} has the form of the Kohn--Sham self-consistency condition
\cite{kohnsham}, which suggested treating it the same way: statistics in place
of density, calibration loss in place of total energy, linear or Anderson
mixing \cite{anderson,pulay} in place of density mixing, and codebook-assignment
churn in place of density convergence. We implemented that loop, with a guard
that rejects any iterate raising the calibration loss and halves the mixing
parameter.

This section reports that the analogy is wrong in the specific way that matters,
and that the repair we designed for it also fails.

\subsection{Why one pass suffices: acyclicity}
\label{sec:acyclic}

The statistics $H^{xx}_\ell$ depend only on layers $1,\dots,\ell-1$. Quantizing
in topological order while propagating the quantized activations forward
therefore gives every layer exactly the statistics it will see in the deployed
model. The dependency graph is acyclic, the fixed point of \eqref{eq:fp} is
reached in \textbf{one sweep}, and there is nothing for an iteration to do.

This is not merely a theoretical point. Two published methods already exploit
it: GPTAQ's asymmetric calibration \cite{gptaq} and CoreQ's mismatch correction
\cite{coreq} both target accumulated error from previously quantized layers, and
both are closed-form and single-pass. We should have derived the acyclicity
argument before building an iteration, and we state it here because it is the
reason the iteration cannot be the contribution.

\paragraph{Measured.} Table~\ref{tab:phase0a} compares the one-pass sequential
arm against our iterative loop on the 2-D codec, the codec most sensitive to the
statistics. The prediction was redundancy; the measurement is stronger. The
one-pass arm beats the iteration by $0.004228$ NLL and beats one-shot
calibration by $0.002778$.

\begin{table}[ht]
\centering\small
\caption{Forward-only self-consistency. Qwen3.5-0.8B, 2-D codec, all 18 DeltaNet
input projections, 32{,}768 WikiText-2 validation targets.}
\label{tab:phase0a}
\begin{tabular}{lrrl}\toprule
Calibration & NLL & $\Delta$NLL & 95\% CI\\\midrule
BF16 reference & 3.233953 & --- & ---\\
one-shot (BF16 activations) & 3.245136 & $+0.011184$ & $[+0.009132,+0.013299]$\\
\textbf{sequential, one pass} & \textbf{3.242358} & $\mathbf{+0.008406}$ & $[+0.006635,+0.010215]$\\
iterative SCF loop & 3.246586 & $+0.012633$ & $[+0.010697,+0.014610]$\\\bottomrule
\end{tabular}
\end{table}

The mechanism is clear in hindsight. Sequential calibration uses statistics that
are \emph{correct}, not approximate. Our loop mixes a \emph{global} statistics
vector toward $H[Q]$ with $\alpha<1$; the mixed statistic corresponds to no
actual model state -- it interpolates between two different networks -- and all
eighteen layers are refitted against it simultaneously. Damping is what makes
SCF stable in density functional theory; here it is what makes it wrong, because
the exact answer is available in one sweep and damping moves away from it.

\paragraph{But solving the forward fixed point correctly does not help much
either.} Table~\ref{tab:phase0a} is a 32{,}768-target screening slice. Repeating
the comparison on the complete validation stream removes the advantage:
sequential calibration changes the Lloyd scalar code by $+0.000022$, leaves
decoupled polar \emph{bit-identical} (that codec ignores activation statistics
entirely), improves repaired polar by $-0.000596$, and \emph{worsens} the 2-D
codec by $+0.000607$ (Table~\ref{tab:main}). The effects are small and change
sign across codecs. The iterative loop behaves the same way across calibration
sizes: on the full stream it worsened the 2-D codec by $+0.002720$ with sixteen
shared calibration windows and improved it by $-0.000585$ with sixty-four
windows and a disjoint guard set.

The explanation is in the residual itself. We replace one tensor in each of the
eighteen DeltaNet layers of this 24-layer model -- $15.1\%$ of its parameters --
and the forward statistics move by only $0.87\%$, so there is very little for any
forward-only correction to recover. Sequential propagation is the
\emph{correct} way to solve that fixed point -- the acyclicity argument is
unaffected -- but at this scope the fixed point is nearly where one-shot
calibration already sits. We would expect the correction to matter more under
whole-model quantization, which we did not run.

We also note the methodological point: the screening slice favoured sequential
calibration by $0.002778$ NLL and the full stream did not reproduce it. Reduced
evaluation slices in this project have repeatedly produced effects that vanish
at full scale, and we treat the screening result as superseded.

\subsection{The bidirectional reformulation}

The fixed point becomes non-trivial exactly when the per-layer objective depends
on layers \emph{downstream}. Expanding the calibration loss to second order in
each layer's output perturbation $\Delta y_\ell=\Delta W_\ell x_\ell$, with the
Gauss--Newton/Fisher approximation to the output Hessian,
\begin{equation}
 \text{cost}_\ell=\sum_j S_{\ell,j}\,\Delta w_{\ell,j}^\top H^{xx}_\ell \Delta w_{\ell,j},
 \qquad
 S_{\ell,j}=\mathbb E\!\left[\left(\frac{\partial\mathcal L}{\partial y_{\ell,j}}\right)^{\!2}\right].
\end{equation}
Here $H^{xx}_\ell$ depends on upstream layers and $S_{\ell,j}$ on downstream
ones, because the gradient at layer $\ell$ is backpropagated through every later
layer and those are quantized too. The dependency graph is now \textbf{cyclic}:
no topological order exists, so no single sweep in either direction satisfies
both conditions. This is the Kohn--Sham situation proper -- the potential
depends on the density everywhere, not only upstream -- and it is the part that
\cite{gptaq,coreq} do not address, because a layer-local output-matching target
assumes all output errors matter equally, which is what $S_\ell$ denies.

\paragraph{The coupling is large.} We measured the relative residual of each
statistic when the eighteen projections are quantized. The forward residual is
$0.87\%$ pooled and grows monotonically with depth ($0.336\%$ at layer 1 to
$1.582\%$ at layer 22), consistent with accumulated upstream error. The backward
residual is $7.33\%$ pooled, $8.41\times$ larger, and is not monotone in depth.
Layer 0 is the clean illustration: its forward residual is \emph{identically
zero}, because its input is upstream of every quantized layer, while its
backward residual is $6.95\%$ because everything downstream of it is quantized.

So the coupling the reformulation identifies is real and an order of magnitude
larger than the one the literature corrects. That made the next test worth
running, and it is where the direction dies.

\subsection{Fisher weighting does not rank layers}
\label{sec:fisherfail}

If $S$-weighted error is the right local surrogate for end-to-end damage, it
should predict which layers are costly to quantize. We quantized one layer at a
time, all others BF16, and correlated four predictors against the measured
validation damage over the eighteen layers.

\begin{table}[ht]
\centering\small
\caption{Predicting single-layer quantization damage. 65{,}536 validation
targets; 10 of 18 layers have a $\Delta$NLL interval excluding zero.}
\label{tab:fisher}
\begin{tabular}{lrrr}\toprule
Predictor & Spearman & Pearson & Pearson without L0\\\midrule
Weight MSE & $-0.218$ & $-0.202$ & $+0.072$\\
Diagonal activation-weighted & $+0.020$ & $-0.061$ & $+0.146$\\
Exact output error (forward only) & $+0.005$ & $-0.057$ & $+0.145$\\
Fisher (forward $\times$ backward) & $+0.015$ & $+0.898$ & $+0.181$\\\bottomrule
\end{tabular}
\end{table}

The Fisher Pearson coefficient of $0.898$ is an artifact of a single leverage
point. Excluding layer 0 it collapses to $0.181$, and its Spearman coefficient
is $-0.169$. The disagreement between a high Pearson and a null Spearman is the
diagnostic: no predictor, Fisher included, has rank skill across the other
seventeen layers.

Two limits on this test's power should be stated. Eight of eighteen layers have
damage intervals covering zero, so the comparison among low-damage layers is
weak. And three layers have \emph{significantly negative} damage -- quantizing
them improves validation NLL -- which no unsigned error surrogate can predict at
all. That caps the achievable correlation independently of the metric.

\subsection{The one real finding, and why it is not exploitable}

Layer 0 is simultaneously the most damaging layer to quantize, the layer Fisher
ranks highest, and the layer weight MSE ranks \emph{lowest} -- that is, the
criterion used throughout the earlier phases of this work identifies the single
most damaging tensor in the model as the safest one. Layer 0 alone accounts for
$58.2\%$ of the summed single-layer damage. This is a crisp mechanism for the
observation, recurring throughout this project, that weight MSE fails to predict
task loss.

It does not convert into a method. Table~\ref{tab:protect} asks which single
tensor is worth leaving in BF16. Fisher's choice recovers $57.7\%$ of the damage
and weight MSE's choice recovers $3.1\%$, which looks decisive until the right
control is added. Leaving one tensor unquantized costs $1.167\times$ the
all-quantized payload, so the honest comparison is a larger codebook applied
uniformly at the same budget.

\begin{table}[ht]
\centering\small
\caption{Protecting one tensor versus spending the same bytes uniformly.
Complete validation stream, 261{,}248 targets.}
\label{tab:protect}
\begin{tabular}{lrlrr}\toprule
Arm & $\Delta$NLL & 95\% CI & Recovered & Payload\\\midrule
All 18 quantized & $+0.008393$ & $[+0.007720,+0.009069]$ & --- & $1.000\times$\\
Protect L0 (Fisher's pick) & $+0.003716$ & $[+0.003197,+0.004245]$ & $55.7\%$ & $1.167\times$\\
Protect L20 (weight MSE's pick) & $+0.008061$ & $[+0.007401,+0.008729]$ & $4.0\%$ & $1.167\times$\\
Protect L10 (arbitrary) & $+0.008382$ & $[+0.007733,+0.009025]$ & $0.1\%$ & $1.167\times$\\
\textbf{All 18, 4.5-bit codebook} & $+0.004872$ & $[+0.004348,+0.005402]$ & $41.9\%$ & $1.125\times$\\
\textbf{All 18, 5-bit codebook} & $\mathbf{+0.002950}$ & $[+0.002555,+0.003346]$ & $\mathbf{64.9\%}$ & $1.250\times$\\\bottomrule
\end{tabular}
\end{table}

Interpolating the uniform-bits frontier to the same $1.167\times$ payload gives
$49.6\%$ recovery. Protecting the Fisher-identified layer is therefore worth
$+6.1$ percentage points, against a confidence half-width of $6.2$ points on the
protected arm alone, before accounting for the interpolation's own uncertainty.
\textbf{The advantage is at best marginal and we do not regard it as
established: outlier protection does not clearly beat spending the same bytes
uniformly.}

\subsection{Summary of Part II}

The forward-only self-consistency loop is redundant by acyclicity, already
addressed in closed form by published one-pass methods, and measurably worse
than the one-pass baseline. The bidirectional reformulation identifies a
genuinely cyclic coupling that is $8.41\times$ larger than the one the
literature corrects, but the metric it implies has no rank skill across layers,
and its single correct call yields no advantage over uniform bit allocation at
matched bytes. We do not recommend pursuing the direction further.


\section{What Survives}
\label{sec:survives}
Neither hypothesis survived. Three findings did, none of them the thing we set
out to show. We state each with the evidence that supports it and the limit of
what it licenses.

\paragraph{1. Uniformly spending a few more bits beat every layer-selective
scheme we tried.} At $1.25\times$ the four-bit payload, a 1024-point 2-D
codebook applied to every layer recovers $64.9\%$ of the quantization damage
(Table~\ref{tab:protect}). No selective scheme we tested -- codec selection per
layer, mixed polar/GPTQ codecs, greedy cumulative selection, Fisher-guided
protection -- delivered more per byte. The negative half of this is equally
useful: it is the control that every layer-selective proposal in this project
should have been run against, and was not.

\paragraph{2. The most damaging layer to quantize is the one weight MSE calls
safest.} Layer 0 ranks first of eighteen in measured damage, first by the Fisher
metric, and \emph{last} by weight MSE -- that is, the criterion used throughout
the earlier phases of this work identifies the single most damaging tensor in
the model as the safest one. It alone accounts for $58.2\%$ of summed
single-layer damage, and protecting it recovers $55.7\%$ of the damage from
quantizing all eighteen.

This is a concrete mechanism for an observation this project made repeatedly
without explaining: that weight-error rankings do not predict task-loss
rankings. It is not a method. As Section~\ref{sec:fisherfail} reports, the
Fisher metric that gets layer 0 right has no rank skill across the other
seventeen layers, and exploiting the one correct call does not clearly beat
finding 1.

\paragraph{3. Paired 2-D vector quantization reaches parity with a well-tuned
scalar code, and the polar parameterization does not.} At four index bits per
weight, the activation-weighted 256-point 2-D codebook gives the best point
estimate of anything we tested. Its advantage over the Lloyd-fitted scalar code
is $+0.000846$ with interval $[-0.000169,+0.001822]$ -- not significant -- while
its advantage over the polar code is $+0.008210$ $[+0.007205,+0.009186]$. The
honest summary is parity with a well-tuned scalar baseline, plus a decode path
(a 256-entry lookup) that is friendlier to a fused kernel than polar's per-pair
sine and cosine.

\paragraph{What did not survive.} We had expected to report sequential
calibration as a free improvement, on the strength of a 32{,}768-target screen
where it beat one-shot calibration by $0.002778$ NLL. It does not replicate on
the complete validation stream, where its effect is at most $\pm0.0006$ NLL and
changes sign across codecs (Section~\ref{sec:acyclic}). We report it here as a
non-finding rather than omit it, because the screen-to-full-stream reversal is
itself the kind of result this paper is about.


\section{Discussion: Control Strength as the Recurring Failure Mode}
\label{sec:discussion}

Both hypotheses in this paper produced apparently significant, well-powered,
reproducible positive results that evaporated when a single control was
strengthened. The pattern is worth stating explicitly because it was not obvious
from inside either project.

\paragraph{The control absorbed the tuning.} In Part~I the polar code had an
implicitly adapted codebook -- annular sectors fitted by a radius grid -- while
the control had a uniform grid. In Part~II the layer-protection scheme was
compared against a different \emph{selection criterion} rather than against the
obvious alternative use of the same bytes. In both cases the proposed method and
the control differed in two ways at once -- the mechanism under test, and the
amount of fitting applied -- and the second difference explained the result.

\paragraph{Diagnostics that would have caught it earlier.} Three cheap checks
found both errors, and none required new infrastructure.
\begin{enumerate}
\item \textbf{Give the control the method's own advantage.} Fitting a Lloyd--Max
codebook for the scalar code costs 16 floats per tensor and reversed the
headline result.
\item \textbf{Measure the ceiling before optimizing inside it.} A 256-point
$k$-means codebook is a few lines and bounds every paired 8-bit scheme. It
showed the polar family capped and the pairing family open, which no amount of
tuning within the polar family could have revealed.
\item \textbf{Match the budget, not the method.} When an arm spends more bytes,
interpolate the uniform-spending frontier to the same budget. This converted
``protecting the right layer recovers $57.7\%$'' into ``$+5.1$ points against an
$11.8$-point noise band''.
\end{enumerate}

\paragraph{A statistical tell.} In Section~\ref{sec:fisherfail} a Pearson
correlation of $0.898$ coexisted with a Spearman correlation of $0.015$. That
disagreement is diagnostic of a single leverage point and should be treated as
an alarm rather than a result; reporting the rank statistic alongside the linear
one would have prevented the overclaim.

\paragraph{On reproducing a prior number.} Our reproduction of the earlier
draft's polar result at its own calibration size gave $+0.011358$ against a
published $+0.009333$, with three seeds spanning $[+0.010515,+0.011358]$. A
published value outside the spread of its own reproduction is evidence about the
selector's variance rather than about the method, and small-calibration
protocols should report a seed sweep for that reason.

\section{Limitations}

One checkpoint, one corpus, one projection family, one evaluation objective.
Earlier phases of this project established that all four axes can change a codec
ranking, so these results are a controlled comparison on a fixed protocol rather
than a general claim. Specifically:

\begin{itemize}
\item \textbf{Scope.} We quantize one tensor in each of the eighteen DeltaNet
layers of a 24-layer model: $113{,}246{,}208$ of 752M parameters, or $15.1\%$.
That is why the forward statistics move by only $0.87\%$ and why sequential
calibration has so little to recover. Whole-model quantization would enlarge
that gap and could change the Part~II conclusions about the forward half, though
not the acyclicity argument itself.
\item \textbf{No speedup.} Task-quality evaluation decodes candidate weights to
BF16 and uses ordinary BF16 matrix multiplies. Byte counts are storage results.
Our packed polar kernels were $3$--$4.5\times$ slower than resident BF16; a
lookup-table decoder removes the trigonometric cost that caused this, but we
have not built or benchmarked that kernel.
\item \textbf{Power.} The single-layer diagnostic in
Section~\ref{sec:fisherfail} leaves eight of eighteen layers with damage
intervals covering zero, so its null result among low-damage layers is weakly
powered. Three layers have significantly \emph{negative} damage, which caps the
correlation any unsigned surrogate can achieve.
\item \textbf{Diagonal statistics.} We mix and weight by diagonal second
moments. Full or block Hessians are a closer analogue of the DFT density and may
behave differently, though the acyclicity argument of
Section~\ref{sec:acyclic} is independent of that choice.
\item \textbf{Controls.} GPTQ- and AWQ-style baselines are local
implementations, not official backend results, and we did not implement GPTAQ or
CoreQ.
\end{itemize}

\section{Conclusion}

Neither hypothesis survived. Polar magnitude--phase weight coding loses to a
real 4-bit code once the scalar code receives the same per-tensor codebook
fitting, and it is structurally capped several percent above that control no
matter how its assignment rule, radius grid, or bit split are tuned.
Self-consistent calibration is the wrong frame for the forward-only problem,
whose dependency graph is acyclic and which one sequential pass solves exactly;
the bidirectional reformulation identifies a coupling $8.41\times$ larger than
the one the literature corrects, but the metric it implies neither ranks layers
nor yields an allocation that beats spending the same bytes uniformly.

What we would keep is smaller than what we set out to show. Uniformly
spending a few more bits beat every layer-selective scheme we tried. The most
damaging layer to quantize is the one weight MSE calls safest, which is a
concrete mechanism for a failure mode this project observed for months without
explaining. And both hypotheses failed in the same way, against controls that
had been given less tuning than the methods they were testing -- a pattern we
would now check for first rather than last.

\appendix
\section{Operator Pilot}
\label{sec:operator}

We retain the earlier sampled-operator study as motivation only. A recurrent
weight is applied repeatedly, so accurate one-step reconstruction need not
preserve long-horizon behavior; in a complex diagonal transition, eigenvalue
magnitude controls decay and phase controls oscillation.

\begin{proposition}[Phase drift for a fixed radius]
If $\lambda=\rho e^{i\phi}$ and $\widehat\lambda=\rho e^{i\widehat\phi}$, then
for integer $T\geq0$,
$|\lambda^T-\widehat\lambda^T| = 2\rho^T|\sin(T(\phi-\widehat\phi)/2)|$.
\end{proposition}
\begin{proof}
Factor out $\rho^Te^{iT(\phi+\widehat\phi)/2}$ and apply $e^{iu}-e^{-iu}=2i\sin u$.
\end{proof}

This bounds modal drift by $\rho^T\min(2,T|\phi-\widehat\phi|)$ and motivated
calibrating against a multi-step impulse response rather than one-step
coefficients. We sampled three systems with $N=64$ complex diagonal modes,
calibration horizon $128$, and evaluation horizons to $512$, comparing
phase-grid and Cartesian decoders under weight and memory objectives at matched
index budgets with binary export. Per-system numbers are in
\texttt{memory\_results.tex}. Three seeds support descriptive statistics only.
Because Part~I shows the phase parameterization is not competitive as a weight
code, we did not develop this pilot further.

\begingroup
\small
\begin{thebibliography}{99}
\setlength{\itemsep}{2pt}
\bibitem{phaseweights} \emph{Artificial neural networks using complex numbers and phase encoded weights}. Applied Optics, 2010. \url{https://opg.optica.org/ao/abstract.cfm?uri=ao-49-10-b71}.
\bibitem{deepcomplex} C. Trabelsi et al. \emph{Deep Complex Networks}. 2017. \url{https://arxiv.org/abs/1705.09792}.
\bibitem{pct} L. Hioki. \emph{Complex-Valued Phase-Coherent Transformer}. 2026. \url{https://arxiv.org/abs/2605.10123}.
\bibitem{urnn} M. Arjovsky, A. Shah, and Y. Bengio. \emph{Unitary Evolution Recurrent Neural Networks}. ICML, 2016. \url{https://arxiv.org/abs/1511.06464}.
\bibitem{s4d} A. Gu, A. Gupta, K. Goel, and C. R\'e. \emph{On the Parameterization and Initialization of Diagonal State Space Models}. 2022. \url{https://arxiv.org/abs/2206.11893}.
\bibitem{qs4d} S. Siegel et al. \emph{QS4D: Quantization-aware training for efficient hardware deployment of structured state-space sequential models}. 2025 preprint; journal version 2026. \url{https://arxiv.org/abs/2507.06079}.
\bibitem{quamba2} H.-Y. Chiang et al. \emph{Quamba2: A Robust and Scalable Post-training Quantization Framework for Selective State Space Models}. 2025. \url{https://arxiv.org/abs/2503.22879}.
\bibitem{gdn} S. Yang, J. Kautz, and A. Hatamizadeh. \emph{Gated Delta Networks: Improving Mamba2 with Delta Rule}. ICLR, 2025. \url{https://arxiv.org/abs/2412.06464}.
\bibitem{gdnquant} S. Kozyrev and D. Maiboroda. \emph{Why Gated DeltaNet Survives 4-Bit Quantization: NVFP4 W4A4 for the Recurrent Half of a Hybrid 27B LLM}. 2026 preprint. \url{https://arxiv.org/abs/2609.04098}.
\bibitem{gptq} E. Frantar, S. Ashkboos, T. Hoefler, and D. Alistarh. \emph{GPTQ: Accurate Post-Training Quantization for Generative Pre-trained Transformers}. ICLR, 2023. \url{https://arxiv.org/abs/2210.17323}.
\bibitem{awq} J. Lin et al. \emph{AWQ: Activation-aware Weight Quantization for On-Device LLM Compression and Acceleration}. MLSys, 2024. \url{https://arxiv.org/abs/2306.00978}.
\bibitem{gptaq} \emph{GPTAQ: Efficient Finetuning-Free Quantization for Asymmetric Calibration}. 2025. \url{https://arxiv.org/abs/2504.02692}.
\bibitem{coreq} \emph{CoreQ: Learning-Free Mismatch Correction and Successive Rounding for Quantization}. 2026. \url{https://arxiv.org/abs/2602.05902}.
\bibitem{kohnsham} W. Kohn and L. J. Sham. \emph{Self-Consistent Equations Including Exchange and Correlation Effects}. Phys. Rev., 1965.
\bibitem{anderson} D. G. Anderson. \emph{Iterative Procedures for Nonlinear Integral Equations}. JACM, 1965.
\bibitem{pulay} P. Pulay. \emph{Convergence acceleration of iterative sequences: the case of SCF iteration}. Chem. Phys. Lett., 1980.
\end{thebibliography}
\endgroup
\end{document}
